%==============================================================================
\section{Backend: Django + Channels}
%==============================================================================

\subsection{Why ASGI and Channels}
A classic Django request/response cycle cannot hold a long-lived, bidirectional
connection open, which is exactly what streaming ECG needs. The project therefore
runs on the asynchronous \textbf{ASGI} stack: \textbf{Django Channels} adds
WebSocket support and \textbf{Daphne} serves the application. The same server
process handles ordinary HTTP (to serve the dashboard page) and WebSockets (to
carry ECG data).

The protocol routing is defined in \texttt{holter\_monitor/asgi.py}: HTTP goes to
the standard Django application, while WebSocket connections are routed through
the ECG app's URL patterns.

\begin{lstlisting}[style=py,caption={ASGI routing -- HTTP and WebSocket on one server.}]
application = ProtocolTypeRouter({
    "http": get_asgi_application(),
    "websocket": AuthMiddlewareStack(
        URLRouter(websocket_urlpatterns)   # from ecg/routing.py
    ),
})
\end{lstlisting}

The WebSocket route (\texttt{ecg/routing.py}) maps the path \texttt{ws/ecg/} to the
\texttt{ECGConsumer}:
\begin{lstlisting}[style=py,numbers=none]
websocket_urlpatterns = [
    re_path(r'ws/ecg/$', ECGConsumer.as_asgi()),
]
\end{lstlisting}

\subsection{The ECG consumer}
\texttt{ECGConsumer} (\texttt{ecg/consumers.py}) is the heart of the backend. Both
the ESP32 \emph{and} every dashboard browser connect to the same
\texttt{ws/ecg/} endpoint, and all of them are placed into a single Channels
\emph{group} called \texttt{ecg\_group}. This group is what lets one producer (the
ESP32) fan its data out to many consumers (the dashboards):

\begin{itemize}[noitemsep]
  \item \textbf{\texttt{connect()}} -- accepts the socket and adds the channel to
        \texttt{ecg\_group}.
  \item \textbf{\texttt{receive()}} -- parses an incoming \texttt{\{"ecg":[...]\}}
        message, runs it through \texttt{filter\_ecg\_signal()}, and broadcasts the
        cleaned array to the whole group via \texttt{group\_send}.
  \item \textbf{\texttt{send\_ecg()}} -- the group handler; serialises the cleaned
        array as \texttt{\{"filtered":[...]\}} and pushes it down each socket.
  \item \textbf{\texttt{disconnect()}} -- removes the channel from the group.
\end{itemize}

\begin{lstlisting}[style=py,caption={Receiving raw ECG, filtering it, and broadcasting to the group.}]
class ECGConsumer(AsyncWebsocketConsumer):
    async def connect(self):
        await self.channel_layer.group_add("ecg_group", self.channel_name)
        await self.accept()

    async def receive(self, text_data):
        data = json.loads(text_data)
        raw  = data.get("ecg", [])
        filtered = filter_ecg_signal(raw)          # DSP pipeline (Section 4)
        await self.channel_layer.group_send(
            "ecg_group",
            {"type": "send_ecg", "data": filtered}, # -> send_ecg()
        )

    async def send_ecg(self, event):
        await self.send(text_data=json.dumps({"filtered": event["data"]}))

    async def disconnect(self, close_code):
        await self.channel_layer.group_discard("ecg_group", self.channel_name)
\end{lstlisting}

\subsection{Channel layer}
Groups are backed by a \emph{channel layer}. The project uses the in-memory layer,
configured in \texttt{settings.py}:
\begin{lstlisting}[style=py,numbers=none]
CHANNEL_LAYERS = {
    "default": {"BACKEND": "channels.layers.InMemoryChannelLayer"}
}
\end{lstlisting}
The in-memory layer is simple and dependency-free, which is ideal for a single
development server. It does not, however, share state across multiple processes,
so a production deployment would swap in the Redis channel layer
(see Section~\ref{sec:limitations}).

\subsection{Serving the dashboard page}
The dashboard itself is plain HTTP. A single view renders the template, wired to
the site root in \texttt{holter\_monitor/urls.py}:
\begin{lstlisting}[style=py,numbers=none]
# ecg/views.py
def ecg_view(request):
    return render(request, "ecg.html")

# holter_monitor/urls.py
urlpatterns = [ path('', ecg_view) ]
\end{lstlisting}

\subsection{Data model}
The current build streams in real time and does \emph{not} persist samples:
\texttt{ecg/models.py} defines no models and there are no migrations beyond
Django's built-ins. Adding stored, per-patient recordings is the most obvious
extension and is listed in Section~\ref{sec:limitations}.
